\section{Renewal energies and admissible suffix trees}\label{sec:v6-suffix}
The main combinatorial fact is a strict energy inequality, not strict
inequality of cylinder masses. The distinction matters at vertices
having a unique predecessor.

\begin{lemma}[Strict suffix domination]\label{lem:v6-suffix-energy}
If $u$ is nonempty and $uv$ is admissible, then
\begin{equation}
 V_w(v)-V_w(uv)\ge
 (1-r_+^\tau)r_v^\tau\sum_{k<|u|}w_k>0,
 \qquad a_w(uv)<a_w(v).                                  \label{eq:v6-strict}
\end{equation}
There are constants $a_0>0$ and $\rho<1$ such that each admissible child
satisfies
\begin{equation}
 a_0a_w(v)\le a_w(vi)\le\rho a_w(v),\qquad
 \sum_{i:\,vi\text{ admissible}}a_w(vi)\le\rho a_w(v).
                                                        \label{eq:v6-child}
\end{equation}
The same inequalities hold at the root after decreasing $a_0$.
\end{lemma}
\begin{proof}
Write $m=|u|$, $n=|v|$ and suppress the subscript $w$. Every suffix of
$v$ has scale at least $r_v$. Since the weights decrease,
\begin{align*}
 V(v)-\left(\sum_{j<n}w_{m+j}r_{v_{j:}}^\tau+W(m+n)\right)
 &\ge r_v^\tau\left(\sum_{j<n}(w_j-w_{m+j})+W(n)-W(m+n)\right)\\
 &=r_v^\tau\sum_{k<m}w_k.
\end{align*}
The remaining prefix terms in $V(uv)$ are at most
$r_v^\tau\sum_{k<m}w_kr_{u_{k:}}^\tau$, which is at most
$r_+^\tau r_v^\tau\sum_{k<m}w_k$. This proves the first inequality,
also for $v=\varnothing$. Stationarity gives
$\mu[uv]\le\mu\{\omega:\omega_{m:m+n}=v\}=\mu[v]$.
Both masses are positive, so the second inequality follows. Neither
argument completes an inadmissible word by an artificial transition.

For $|v|=n$, appending $i$ multiplies each old spatial scale by a
factor between $D_0^{-1}r_-$ and $r_+$. The old tail $W(n)$ is replaced
by $w_nr_i^\tau+W(n+1)$. The weight assumptions imply
\[
 (1-\overline q)W(n)\le w_n,\qquad
 W(n+1)\ge\underline q W(n).
\]
It follows that
\[
 \ell_0 V(v)\le V(vi)\le\rho V(v),\quad
 \ell_0=\min(D_0^{-\tau}r_-^\tau,\underline q),\quad
 \rho=1-(1-r_+^\tau)(1-\overline q)<1.
\]
A uniform lower bound $\zeta>0$ on admissible conditional cylinder
masses now gives $a_0=\zeta\ell_0$ in \eqref{eq:v6-child}.
The last assertion follows by summing the masses of the children.
Such a $\zeta$ follows directly from \eqref{eq:v6-gibbs} and the
bounded potential, for admissible extensions only.
\end{proof}

\begin{lemma}[Balanced trees and controlled dilation]\label{lem:v6-balanced}
Let $L_N=|\mathcal P_N|$. For all sufficiently large $N$,
$N/b\le L_N\le N$. If $\theta_N$ is the energy of the last split
vertex, then every leaf satisfies
\begin{equation}
        a_0\theta_N\le a_w(v)\le\theta_N.                \label{eq:v6-balanced}
\end{equation}
For every fixed $B\ge1$, with constants allowed to depend on $B$,
\begin{equation}
                   G_{\lceil BN\rceil}(w)\asymp_B G_N(w).
                                                        \label{eq:v6-dilation}
\end{equation}
The complete tree with $L$ leaves has at most $(b+1)(2L-1)$ vertices.
\end{lemma}
\begin{proof}
Maximal split energies never increase, by \eqref{eq:v6-child}.
A present leaf was born with energy at least $a_0$ times an earlier
split energy, which proves the lower bound in
\eqref{eq:v6-balanced}; the upper bound follows from the final
maximal split. Each split adds between zero and $b-1$ leaves.
The split following $\mathcal T_N$ increases the number beyond $N$;
hence $L_N>N-b+1$, which implies the stated coarser bound.

Here are the finiteness details needed in the presence of unary
vertices. A path of $b$ successive vertices all having outdegree one
would repeat a vertex and enter a closed deterministic cycle. This
contradicts primitivity and $\lambda_A>1$. Thus a unary chain has
length at most $b$. After suppressing unary vertices, any finite tree
with $L$ leaves and no internal vertex of degree one has at most
$2L-1$ vertices. Reinserting the unary chains, including an initial
chain, gives the asserted bound. In particular an infinite sequence
of splits with bounded leaf number is impossible.

Choose $\ell$ such that $A^\ell$ has all entries positive. From every
last symbol there are at least two distinct continuations of length
$\ell$. Thus every vertex has at least $2^j$ descendants at depth
$\ell j$. In continuing a greedy tree with threshold $\theta_N$, if
all later leaves have energy less than $a_0^{\ell j+1}\theta_N$,
then every original leaf must have been split through those $\ell j$
levels. This would give at least $2^jL_N$ later leaves. Choosing $j$
so that $2^j> b(B+2)$ shows that the final maximum leaf energy at
budget $\lceil BN\rceil$ is at least a fixed multiple of $\theta_N$.
The balance estimate at that budget bounds its minimum leaf energy
by another fixed multiple of $\theta_N$. Since $L_{\lceil BN\rceil}$
is at least a fixed multiple of $N$, its energy sum is at least
$c_BN\theta_N\ge c_BG_N$. Conversely splitting decreases the sum
by \eqref{eq:v6-child}, so $G_{\lceil BN\rceil}\le G_N$.
Increasing constants handles the finitely many initial trees.
\end{proof}

\begin{lemma}[Observable orbit cylinders]\label{lem:v6-cylinders}
The image $Z_w^f(K_v)$ has squared diameter at most $C V_w(v)$.
If two different sequences have longest common prefix $u$, then
\begin{equation}
        \|Z_w^f(x)-Z_w^f(y)\|^2\ge c V_w(u).             \label{eq:v6-separation}
\end{equation}
All constants depend only on the displayed geometric and observable
hypotheses, not on the word.
\end{lemma}
\begin{proof}
Introduce the auxiliary snowflake orbit distance
$d_\beta(x,y)^2=\sum_{k\ge0}w_k|T^kx-T^ky|^{2\beta}$.
The H\"older estimate bounds $\|Z_w^f(x)-Z_w^f(y)\|^2$ from above
by $H^2d_\beta(x,y)^2$. Multiplying \eqref{eq:v6-delay} at $T^kx,T^ky$
by $w_k$ and summing gives
\[
 c_od_\beta(x,y)^2\le
 \sum_{j=0}^J\sum_{k\ge0}w_k
       |f(T^{k+j}x)-f(T^{k+j}y)|^2
 \le\left(\sum_{j=0}^J\underline q^{-j}\right)
       \|Z_w^f(x)-Z_w^f(y)\|^2.
\]
For points in $K_v$, the first $n$ terms of $d_\beta^2$ are bounded
by $w_kr_{v_{k:}}^{2\beta}$ and its remaining terms by $W(n)$.
If $u$ is their longest common prefix, then at time $|u|$ they lie
in different first-level images, and at time $k<|u|$ their distance
is at least $D_0^{-1}g r_{u_{k:}}$. The first post-prefix term is
$w_{|u|}g^{2\beta}$, and $W(|u|)\le w_{|u|}/(1-\overline q)$.
These estimates prove both assertions, including the empty prefix.
\end{proof}

\begin{lemma}[Unrestricted quantization and suffix realization]
\label{lem:v6-realization}
The quantization errors satisfy $e_M^2(Z_w^f)\asymp G_M(w)$.
Every complete greedy tree is suffix closed, including its leaves.
A tree with $L$ leaves gives a deterministic observer with at most
$(b+1)(2L-1)$ states and average loss at most $C\sum_{v\in\mathcal P}a_w(v)$.
\end{lemma}
\begin{proof}
Taking one orbit representative from each leaf cylinder proves
$e_N^2\le CG_N$. For the reverse direction use the leaf partition
$\mathcal P_{4bM}$, which has at least $4M$ members, apart from
finitely many budgets absorbable in the constants. Around
$Z_w^f(K_v)$ take its open Hilbert tube of radius $\delta\sqrt{V_w(v)}$.
These tubes are disjoint for a sufficiently small fixed $\delta$:
for incomparable leaves the separation in \eqref{eq:v6-separation}
is at their common prefix, and appending symbols only decreases $V_w$.
Any $M$ Hilbert centres, whether or not on the orbit image, occupy at
most $M$ tubes. On every other cylinder the nearest-centre squared
error is at least $\delta^2V_w(v)$. The balance estimate shows that
these unoccupied cylinders carry at least a fixed fraction of the
energy sum. Hence $e_M^2\ge cG_{4bM}\ge c'G_M$ by
\eqref{eq:v6-dilation}. For each of the finitely many omitted budgets
this error is positive: the delay condition and infinitely supported
Gibbs law exclude support on finitely many orbit points.

Suppose an internal word $uv$, with $u$ nonempty, is about to be
split. If its suffix $v$ has not yet been split, some current leaf
$t$ is a prefix of $v$. By \eqref{eq:v6-child} and
\eqref{eq:v6-strict}, $a_w(t)\ge a_w(v)>a_w(uv)$, contradicting
maximality. Thus every proper suffix of an internal word is already
internal. If $v=i_0\cdots i_{n-1}$ is a leaf with $n\ge2$, the word
$i_1\cdots i_{n-2}$ is a suffix of its internal parent, so it has
been split. Its admissible child $i_1\cdots i_{n-1}$ is therefore
present. For $n=1$ the suffix is the root. Iteration proves closure
of the entire tree; it does not require that all symbols follow
all others.

Choose $x_v\in K_v$ for every nonempty tree vertex, and decode that
vertex as $f(x_v)$. At an acquisition store the unique leaf prefix
of the observed sequence. At each nonacquisition delete the first
symbol. Decode the empty word as any fixed $f(x_*)$, and keep it
until the next acquisition. Each update stays within the same tree.
If the acquired word is $v$ of length $n$, the squared loss at age
$k<n$ is at most $H^2r_{v_{k:}}^{2\beta}$; after exhaustion it is
bounded by a constant. The renewal-reward formula therefore bounds
the average by $C\sum_v\mu[v]V_w(v)$. All vertices are charged, with
the count established in Lemma~\ref{lem:v6-balanced}.
\end{proof}

\begin{proof}[Proof of Theorem~\ref{thm:v6-profile}]
The retained posterior-orbit theorem, Theorem~\ref{thm:v4-orbit}, applies
with exact observation and Bayes floor zero. It gives $e_M^2\le\mathcal R_M$
and the stated lower-limit assertion for arbitrary stochastic kernels.
Its full proof is included in both reading editions. Lemma~\ref{lem:v6-realization}
gives the other lower bound. Choose a leaf budget $N$ a fixed sufficiently
small fraction of $M$ so that \eqref{eq:v6-state-count} is at most $M$.
The same lemma and fixed-budget dilation give an upper bound
$CG_N\le C'G_M$. A one-state observer handles the remaining bounded
range of $M$. This proves \eqref{eq:v6-profile-law}.
\end{proof}
